\begin{center}
\textsc{Gauss} an \textsc{Bessel}.\quad Göttingen, 20.~November 1826.
\end{center}

\par\noindent\hspace*{\parindent}.\ .\ .\ .\ .\ .\quad Die Verarbeitung der Materialien zu dem beabsichtigten Werke über meine Messungen kostet mich viele Zeit.
Meine Hauptdreiecke, 33 Punkte befassend, sind zwar längst fertig berechnet, aber die Berechnung der vielen geschnittenen Nebenpunkte, die für die Geographie eines bedeutenden Theils von Norddeutschland wichtig sind, macht viel Arbeit, da jene bisher entweder noch gar nicht oder nur provisorisch berechnet waren.
Mein Verzeichniss enthält jetzt etwa 250 Punkte, alle aus dem nördlichern Theil meiner Messungen.
Da meine Dreiecke durch die kurhessischen mit den bayerischen und württembergischen zusammenhängen, und letztere mir auch von \textsc{Bohnenberger} gefällig mitgetheilt sind (ebenso wie die darmstädtschen von \textsc{Eckhardt}), so ist es mir sehr unangenehm, dass alle meine Bemühungen, die bayerischen zu erhalten,
\par\smallskip\noindent{\footnotesize IX.}\hfill{\footnotesize46}\par
\clearpage
\noindent362\hfill{\small\textbf{BRIEFWECHSEL MIT BESSEL.}}\par\medskip
\noindent bisher vergeblich gewesen sind. \textsc{Schumacher} erzählte mir, dass \textsc{Soldner} ihm gesagt hätte, der Grund, warum die Commission in München meine Bitte um die Mittheilung nicht erfüllt habe, sei, weil man annehme, dass ich über diese Dreiecke Rechnungen anstellen wolle!
Ebenso ist mir's mit den östreichschen gegangen.
Auch General \textsc{von Müffling} sollte doch seine Dreiecke östlich von Seeberg, die Berlin anschliessen und sich bis über Schlesien erstrecken, bekannt machen.

\par\noindent\hspace*{\parindent}Noch viel mehr Verlegenheit macht mir der weit ausgedehntere theoretische Theil, der so vielfach in andere Theile der Mathematik eingreift.
Ich sehe hier kein anderes Mittel, als mehrere grosse Hauptparthien von dem Werke abzutrennen, damit sie selbstständig und in gehöriger Ausführlichkeit entwickelt werden können.
Gewissermaassen habe ich damit schon in meiner Schrift über die Abbildung der Flächen unter Erhaltung der Ähnlichkeit der kleinsten Theile den Anfang gemacht; eine zweite Abhandlung, die ich vor ein paar Monaten der königlichen Societät übergeben habe, und die hoffentlich bald gedruckt werden wird, enthält die Grundsätze und Methoden zur Ausgleichung der Messungen (beiläufig ist daraus indirect auch ersichtlich, wie weit die \textsc{Krayenhoff}schen Messungen von derjenigen Genauigkeit entfernt sind, die man ihnen mit Unrecht beigelegt hat).
Vielleicht werde ich zunächst erst noch eine dritte Abhandlung ausarbeiten, die mancherlei neue Lehrsätze über krumme Flächen, kürzeste Linien, Darstellung krummer Flächen in der Ebene, u.\ s.\ w.\ entwickeln wird.
Hätten alle diese Gegenstände in mein projectirtes Werk aufgenommen werden sollen, so hätte ich entweder manches ungründlich abfertigen oder dem Werk ein sehr buntscheckiges Ansehen geben müssen.\ .\ .\ .\ .\ .\ .
